\chapter{A Language for Computations: Values and Functions}
\label{ch:language}

\chapterorient{In \cref{ch:bigidea} we made a decision: describe a simulation as
a chain of pure functions over immutable tensors, not as overwrites of a block of
memory. That decision earns its keep only if we can write the functions down
\emph{precisely}---precisely enough that the written form \emph{is} the dependency
graph, precisely enough that a critic can check one piece against the source
without running anything. This chapter introduces the small notation we use for
exactly that, and it teaches you to read it from zero. We assume no
functional-programming background at all. By the end you will be able to look at a
line like \lstinline[style=l4style]|axpy :: Scalar -> Tensor[N] -> Tensor[N] -> Tensor[N]|
and say, token by token, what it means---and you will understand the one idea that
trips up every newcomer: a function that takes, or returns, another function.}

\section{Why a language at all}

We could describe every algorithm in this book in ordinary English prose, or in
ordinary mathematics. We will, in fact, lean on both constantly. But neither is
quite enough on its own, and the gap is worth naming, because it tells us exactly
what the small notation of this part is \emph{for}.

Recall the three questions the calculus is built to answer
(\cref{ch:bigidea}): \emph{what operations happen}, \emph{who owns the state},
and \emph{how state evolution is coordinated}. Each is a structural question---a
question about the \emph{shape} of a computation, not about the numbers it
produces. English prose blurs structure: ``then we update the residual'' hides
whether the residual is a fresh value or an overwrite, whether the update reads
any state besides its stated inputs, whether it is a step in a loop or a one-off.
Ordinary mathematics fixes the arithmetic but says nothing about ownership or
coordination at all: the equation $\vr = \vb - \mat{A}\vx$ tells us what $\vr$
\emph{equals}, but not that $\vr$ is a brand-new immutable value, nor that
computing it is a pure function of exactly $\vb$, $\mat{A}$, and $\vx$ and
nothing else.

So we adopt a third notation---a small, deliberately readable \emph{pseudo-language}
borrowed from typed functional programming---that makes the structure
explicit. This chapter does the first of the three questions, the foundation the
other two rest on: \emph{values and functions}. What is a value? What is a
function? How do we write a function's type, supply its arguments, and---the
hard part---pass functions to other functions? \Cref{ch:records} adds records
(bundles of named fields); \cref{ch:tensors} adds shapes; \cref{ch:fold} adds
the loop combinator; \cref{ch:monad} adds the discipline for threaded state. Each
builds on this one.

\begin{notationbox}
The pseudo-language is a language you \emph{read}, not one you run. It always
appears in a shaded, monospaced box. It is not a burn API, not a Rust sketch, not
something a compiler accepts; it is a precise notation for talking about
algorithms, the way a physicist reads a Feynman diagram. The notation primer at
the front of the book sketched it in one page; this chapter is the careful
version. When you see a shaded box you are reading code; when you see an inline
formula like $\mat{K}\vx=\vb$ you are reading mathematics. The two never mix in
one expression.
\end{notationbox}

\section{Values: things you can name but never change}

The most foundational idea in the whole notation is also the simplest, and you
already met it in \cref{ch:bigidea}: a \emph{value} is a thing you can name but
never change.

The number $7$ is a value. So is the number $3$. When you write $3 + 4$ you do not
\emph{turn} the $3$ into a $7$; you compute a \emph{new} value, $7$, and the $3$
and the $4$ sit there, untouched, exactly as before. Nobody is tempted to ask
``but what is the $3$ \emph{now}, after the addition?''---the question is
nonsense. A number has no ``now.'' It simply is what it is.

The single rule of our notation is that \emph{everything} behaves this way. A
tensor is a value, like $7$. A record (a bundle of named fields, \cref{ch:records})
is a value. A function---we will see---is a value. None of them can be changed
once it exists. You can compute new values \emph{from} them, name those new
values, and pass them around; you can never reach inside an existing value and
overwrite a part of it.

\begin{definition}[Value and immutability]
A \emph{value} is a piece of data---a number, a tensor, a record, a function---that,
once it exists, never changes. Computing with values means producing \emph{new}
values from old ones. A name, once bound to a value, refers to that same
unchanging value for as long as the name is in scope. We call this discipline
\emph{immutability}.
\end{definition}

\subsection{The contrast: a variable as a mutable box}

If you have programmed in C, Python, or most mainstream languages, you carry a
\emph{different} mental model of what a name means, and it is worth dragging into
the light because it will fight you here. In those languages a variable is a
\emph{box}: a labeled slot in memory that holds a value \emph{now} but can be
made to hold a different value \emph{later}.

\begin{lstlisting}[style=l4style, language=]
y = 5      # the box labeled y now holds 5
y = y + 1  # overwrite: the box labeled y now holds 6
\end{lstlisting}

\noindent After the second line, the box \texttt{y} holds $6$; the $5$ is gone.
The \emph{name} \texttt{y} survived but what it points to was destroyed and
replaced. This is the in-place-mutation picture of \cref{ch:bigidea}, now at the
level of a single variable rather than a whole array: one box, scribbled over.

Our notation has no such boxes. When we write

\begin{lstlisting}[style=l4style]
let y = 5 in
let z = y + 1 in
...
\end{lstlisting}

\noindent the name \texttt{y} is bound, once and forever (within its scope), to
the value $5$. We do not---cannot---``update \texttt{y} to $6$.'' Instead we
introduce a \emph{new} name, \texttt{z}, for the new value $6$. There are two
names and two values; nothing was overwritten. The \code{let} keyword reads
``\emph{let} this name stand for this value, in what follows.''

\begin{figure}[t]
\centering
\begin{tikzpicture}[node distance=8mm]
  % ============ LEFT: mutable box ============
  \begin{scope}
    \node[font=\bfseries\small, simRust] at (0,2.5) {A mutable box (a variable)};
    \node[draw=simRust, fill=simBgRust, thick, rounded corners=2pt,
          minimum width=24mm, minimum height=11mm] (box) at (0,1.1) {};
    \node[font=\ttfamily\small, simInk] at (-0.7,1.1) {y};
    \node[font=\small, simInk] at (0.45,1.1) {$\leftarrow 5$};
    % overwrite loop
    \draw[backflow] (box.north west) to[out=135,in=45,looseness=4.5]
      node[above,font=\scriptsize]{\texttt{y := 6}} (box.north east);
    \node[align=center, font=\scriptsize\itshape, simGray, text width=40mm]
      at (0,-0.15) {one slot; the old contents are destroyed and replaced.
      ``What is \texttt{y} now?'' is a real question.};
  \end{scope}
  % ============ RIGHT: immutable values ============
  \begin{scope}[shift={(7.2,0)}]
    \node[font=\bfseries\small, simBlue] at (0,2.5) {Immutable values (names)};
    \node[draw=simBlue, fill=simBgBlue, thick, rounded corners=2pt,
          minimum width=12mm, minimum height=10mm] (v5) at (-0.9,1.1) {$5$};
    \node[draw=simBlue, fill=simBgBlue, thick, rounded corners=2pt,
          minimum width=12mm, minimum height=10mm] (v6) at (1.1,1.1) {$6$};
    \node[font=\ttfamily\scriptsize, simBlue, below=0.5mm of v5] {y};
    \node[font=\ttfamily\scriptsize, simBlue, below=0.5mm of v6] {z};
    \draw[flow] (v5) -- node[above,font=\scriptsize]{$+1$} (v6);
    \node[align=center, font=\scriptsize\itshape, simGray, text width=42mm]
      at (0.1,-0.15) {two distinct values, two names; \texttt{y} still
      means $5$. ``What is \texttt{y} now?'' is a non-question.};
  \end{scope}
\end{tikzpicture}
\caption[A mutable box versus an immutable value]{Two readings of a name.
\emph{Left:} the mainstream model---a variable is a box in memory; \texttt{y := 6}
overwrites it, and the previous contents vanish, so the box's value depends on how
far the program has run. \emph{Right:} our model---a name is bound permanently to
an immutable value; computing $y+1$ produces a \emph{fresh} value $6$, which we
give a \emph{new} name \texttt{z}. The original \texttt{y} still means $5$. The
notation in this book has only the right-hand picture; there are no boxes.}
\label{fig:valuebox}
\end{figure}

\Cref{fig:valuebox} draws the two pictures side by side, the same contrast
\cref{fig:mutvalue} drew for whole arrays, now for a single name. The discipline
is the same at every scale: never overwrite, always produce a fresh value and
name it.

\begin{keyidea}
Immutability is the ground rule of the entire notation. A value---number, tensor,
record, or function---never changes after it is created. A name binds to a value
permanently within its scope. We never overwrite; we compute new values from old
ones and give them new names. Every other convention in this part is built on
top of this one. If you ever find yourself asking ``what does this name hold
\emph{now}?'', you have slipped back into the mutable-box model---stop and
recover the value picture.
\end{keyidea}

\subsection{Why \code{clone()} is meaningless here}

A quick but illuminating consequence. In a mutable language, when you want a
function to work on a copy of your data without disturbing the original, you
\emph{clone} it: you make a second box holding the same contents, so that
overwriting one leaves the other intact. The lattice-Boltzmann reference code in
\cref{ch:bigidea} was littered with \texttt{current.clone()} for exactly this
reason.

In our notation, \code{clone()} is meaningless---there is nothing to clone,
because there is nothing to protect. A value cannot be overwritten, so there is no
way for one use of a value to disturb another use of it. Two names for ``the same
tensor'' is perfectly safe: neither can mutate it, so neither can surprise the
other.

\begin{lstlisting}[style=l4style]
let v = some_tensor in
let a = scale 2 v in   -- uses v
let b = scale 3 v in   -- ALSO uses v; no clone needed
(a, b)                 -- v was never disturbed; it could not be
\end{lstlisting}

\noindent This is not a small convenience. ``When must I copy, to be safe?'' is
one of the central, error-prone questions of mutable programming, and in the
value picture it simply evaporates. The cost we appear to pay---``but surely
making a fresh tensor every step is wasteful!''---is recovered in full by the
implementation, which is free to reuse storage whenever an old value is no longer
referenced (the pitfall in \cref{ch:bigidea}; made precise in \cref{ch:monad}).
We describe the computation with fresh values; the machine optimizes the memory.

\section{Functions and types}

A \emph{function} is a rule that turns inputs into an output. In our notation a
function is written in two parts: a \emph{type signature} that says what kinds of
thing go in and come out, and a \emph{definition} that says how the output is
computed. Here is the smallest possible example---a function that squares a
floating-point number.

\begin{lstlisting}[style=l4style]
square :: Float -> Float       -- the TYPE: takes a Float, gives a Float
square x = x * x               -- the DEFINITION: the output is x times x
\end{lstlisting}

Read the two lines carefully, because their grammar recurs on every function in
the book.

\begin{description}[style=nextline,leftmargin=1.5em]
  \item[The signature \code{square :: Float -> Float}.] The
  double colon \code{::} reads ``\emph{has type}.'' So this line says
  ``\code{square} has type \code{Float -> Float}.'' The arrow \code{->} reads
  ``\emph{maps to}'': \code{Float -> Float} is the type of functions that map a
  \code{Float} to a \code{Float}. The signature is a \emph{contract}: it promises
  that if you hand \code{square} a \code{Float}, you get a \code{Float} back, and
  it forbids handing it anything else.
  \item[The definition \code{square x = x * x}.] Here \code{x}
  names the input (the \emph{parameter}); the right-hand side \lstinline[style=l4style]|x * x|
  is the output, an expression built from the input. Applying the function is
  written by \emph{juxtaposition}---just putting the argument next to the function
  name, with a space: \lstinline[style=l4style]|square 5| evaluates to $25$. No
  parentheses are needed around the argument (though \lstinline[style=l4style]|square (a + b)|
  needs them, to group the argument).
\end{description}

\noindent The types are not decoration. They are the first line of defence
against nonsense: a tool (and a careful reader) can check that every function is
only ever applied to arguments of the type it expects, long before any numbers
are computed. The shapes of tensors, which we develop in \cref{ch:tensors}, are
types in exactly this sense---and they catch the overwhelming majority of
mistakes in tensor code.

\subsection{Pure functions}

The functions in this book obey one further, crucial discipline: they are
\emph{pure}. A pure function's output depends on \emph{nothing but its declared
inputs}, and computing it changes nothing in the outside world. Two consequences,
both load-bearing:

\begin{itemize}[nosep]
  \item \textbf{Same inputs, same output, always.} \lstinline[style=l4style]|square 5|
  is $25$ today, tomorrow, on every machine, no matter what ran before it. A pure
  function has no memory and no mood.
  \item \textbf{No side effects.} A pure function does not secretly read a global
  variable, write to a file, overwrite one of its inputs, or depend on the time
  of day. Everything it uses arrives through its parameters; everything it
  produces leaves through its return value.
\end{itemize}

This is precisely the property that makes the program text equal to the
dependency graph (\cref{ch:bigidea}): if a function can depend only on its
inputs, then drawing an arrow from each input to the function captures
\emph{everything} the function depends on---there are no hidden arrows. Purity is
what we buy with immutability, and the rest of the book spends it.

\begin{pitfall}
The mutable-box habit dies hard, and it shows up as a specific reading error.
When you see \lstinline[style=l4style]|axpy a x y = a * x + y|, the C programmer's
instinct is to ask ``and which of \code{x}, \code{y} did it overwrite?'' The
answer is \emph{neither}. A pure function \emph{returns} its result; it does not
deposit it into one of its arguments. After \lstinline[style=l4style]|let r = axpy a x y|,
the names \code{x} and \code{y} still refer to exactly the tensors they did
before, and \code{r} names the fresh result. If you ever catch yourself hunting
for ``the output argument,'' you have slipped models---there is no output
argument, only a returned value.
\end{pitfall}

\subsection{A function is itself a value}

Here is the move that makes the whole notation expressive, stated plainly now and
exploited heavily later: \emph{a function is a value, just like a number or a
tensor.} It is a thing you can name, pass to another function as an argument, and
return from a function as a result. \code{square} is a value of type
\code{Float -> Float}, exactly as $5$ is a value of type \code{Float}.

This is not how most engineers first learn to think about functions---a function
feels like a \emph{verb}, a thing you \emph{do}, not a \emph{noun}, a thing you
\emph{hold}. But in this notation it is both. \Cref{fig:fnvalue} draws the
cartoon: a function is a little box you can carry around in a bag, hand to
someone, or get handed back. Sections~\ref{sec:currying}
and~\ref{sec:higherorder} are entirely about what this buys us.

\begin{figure}[t]
\centering
\begin{tikzpicture}[node distance=10mm]
  % a "bag" of values, one of which is a function
  \node[font=\bfseries\small, simInk] at (0,2.2) {Values you can name, pass, and return};
  % number value
  \node[draw=simBlue, fill=simBgBlue, thick, rounded corners=2pt,
        minimum width=15mm, minimum height=10mm] (n) at (-3.6,0.7) {$5$};
  \node[font=\scriptsize\itshape, simGray, below=1mm of n] {a number};
  % tensor value
  \node[draw=simBlue, fill=simBgBlue, thick, rounded corners=2pt,
        minimum width=18mm, minimum height=10mm] (t) at (-1.1,0.7)
        {\lstinline[style=l4style]|Tensor[N]|};
  \node[font=\scriptsize\itshape, simGray, below=1mm of t] {a tensor};
  % function value, drawn as a little machine box
  \node[draw=simTeal, fill=simBgGreen, thick, rounded corners=2pt,
        minimum width=30mm, minimum height=12mm] (f) at (2.3,0.7) {};
  \node[font=\ttfamily\scriptsize, simInk] at (1.55,0.95) {square};
  \node[font=\scriptsize, simTeal] at (2.55,0.55)
        {$x \mapsto x{*}x$};
  \node[font=\scriptsize\itshape, simGray, below=1mm of f]
        {a function---also a value};
  % a "carry" arrow suggesting you can pass it around
  \draw[flow] (3.95,0.7) -- ++(1.1,0)
    node[right,font=\scriptsize\itshape,simGray,align=left] {pass it,\\return it,\\name it};
\end{tikzpicture}
\caption[A function is a value you can pass around]{In this notation a function is
a value with the same standing as a number or a tensor: you can bind it to a name,
hand it to another function as an argument, and get one back as a result.
\code{square} is a value of type \code{Float -> Float} the way $5$ is a value of
type \code{Float}. This single idea---functions are first-class values---is what
makes currying (\cref{sec:currying}) and higher-order functions
(\cref{sec:higherorder}) possible.}
\label{fig:fnvalue}
\end{figure}

\section{Currying: one argument at a time}
\label{sec:currying}

Most interesting functions take more than one argument. We write a two-argument
function's type with two arrows:

\begin{lstlisting}[style=l4style]
add :: Float -> Float -> Float
add x y = x + y
\end{lstlisting}

\noindent and apply it by listing the arguments after the name, separated by
spaces: \lstinline[style=l4style]|add 3 4| evaluates to $7$. So far this looks
like the multi-argument functions of any language. But the arrow notation is
hinting at something subtler, and the subtlety pays off enormously, so we take it
slowly.

\subsection{Reading \code{A -> B -> C}}

How should we read \lstinline[style=l4style]|add :: Float -> Float -> Float|? The
honest reading is \emph{one argument at a time}:

\begin{quote}
\emph{``\code{add} takes a \code{Float}, and returns a function that takes a
\code{Float} and returns a \code{Float}.''}
\end{quote}

\noindent That is, the arrow \code{->} groups \emph{to the right}: the type is
secretly
\begin{lstlisting}[style=l4style]
add :: Float -> (Float -> Float)
\end{lstlisting}
\noindent with invisible parentheses around the tail. \code{add} is not really a
two-argument function at all; it is a one-argument function whose \emph{result is
another function}. This way of treating multi-argument functions---as a chain of
one-argument functions---is called \emph{currying}, after the logician Haskell
Curry (whose first name, not coincidentally, also names the language our notation
borrows from).

When you write \lstinline[style=l4style]|add 3 4|, two things happen in
sequence:

\begin{lstlisting}[style=l4style]
add 3       -- supply the first argument: get back a FUNCTION,
            --   "the function that adds 3 to its argument"
(add 3) 4   -- supply the second argument to THAT function: get 7
\end{lstlisting}

\noindent The juxtaposition \lstinline[style=l4style]|add 3 4| is just
\lstinline[style=l4style]|(add 3) 4| with the parentheses dropped---application,
too, groups to the left. So ``calling a two-argument function'' is really
``calling a one-argument function, then calling its one-argument result.'' You
almost never have to think about this when you supply all the arguments at once;
it matters when you supply only \emph{some}.

\subsection{Partial application}

Because \lstinline[style=l4style]|add 3| is itself a perfectly good function value,
we can name it and use it later:

\begin{lstlisting}[style=l4style]
let add3 = add 3 in   -- supply ONE of two arguments; add3 :: Float -> Float
let u = add3 4 in     -- u = 7
let w = add3 10 in     -- w = 13
(u, w)
\end{lstlisting}

\noindent Supplying \emph{some} of a function's arguments and getting back a
function that awaits the rest is called \emph{partial application}. We applied
\code{add} to one argument and \emph{froze} it---baked the $3$ in---producing a
new, more specialized function \code{add3} that adds $3$ to whatever it is later
given. \Cref{fig:currying} draws the freezing.

\begin{figure}[t]
\centering
\begin{tikzpicture}[node distance=9mm]
  % the curried machine, three "slots"
  \node[font=\bfseries\small, simInk] at (0,2.55)
    {Partial application: freeze some arguments now, supply the rest later};
  % stage 0: the full function
  \node[draw=simTeal, fill=simBgGray, thick, rounded corners=2pt,
        minimum width=20mm, minimum height=10mm] (f0) at (-4.3,1.0) {};
  \node[font=\ttfamily\scriptsize] at (-4.3,1.28) {add};
  \node[font=\scriptsize, simTeal] at (-4.3,0.78) {needs $x,y$};
  \node[font=\scriptsize\itshape,simGray,below=0.5mm of f0]
    {\code{Float -> Float -> Float}};
  % supply 3
  \draw[flow] (-3.2,1.0) -- node[above,font=\scriptsize]{supply $x=3$} (-1.7,1.0);
  % stage 1: partially-applied, one slot frozen
  \node[draw=simBlue, fill=simBgBlue, thick, rounded corners=2pt,
        minimum width=22mm, minimum height=10mm] (f1) at (-0.4,1.0) {};
  \node[font=\ttfamily\scriptsize] at (-0.4,1.30) {add3 = add 3};
  \node[font=\scriptsize, simBlue] at (-0.4,0.80) {needs $y$};
  \node[font=\scriptsize\itshape,simGray,below=0.5mm of f1]
    {\code{Float -> Float}};
  % a little "frozen 3" tag
  \node[draw=simRust,fill=simBgRust,rounded corners=1pt,font=\scriptsize,
        inner sep=1.5pt] at (-0.4,2.0) {$x{=}3$ frozen in};
  % supply 4
  \draw[flow] (0.8,1.0) -- node[above,font=\scriptsize]{supply $y=4$} (2.3,1.0);
  % stage 2: a value
  \node[draw=simRust, fill=simBgRust, thick, rounded corners=2pt,
        minimum width=14mm, minimum height=10mm] (r) at (3.4,1.0) {$7$};
  \node[font=\scriptsize\itshape,simGray,below=0.5mm of r] {a \code{Float}};
\end{tikzpicture}
\caption[Currying and partial application]{A curried function is consumed one
argument at a time. \code{add} needs two \code{Float}s; supplying the first
($x=3$) does not produce a number---it produces a \emph{new function} \code{add3},
of type \code{Float -> Float}, with the $3$ \emph{frozen in}. Supplying the
second ($y=4$) to \emph{that} function finally produces the value $7$. Freezing
some arguments now and supplying the rest later is \emph{partial application}; it
is the mechanism behind ``build once with configuration baked in, apply many
times.''}
\label{fig:currying}
\end{figure}

\subsection{Why partial application matters}

Partial application looks like a curiosity, but it is one of the most useful
patterns in the entire book, and it is worth seeing why \emph{now}, even before we
can spell the example fully. Over and over, an algorithm has a piece that is
\emph{configured once} and then \emph{applied many times}: a linear operator
assembled from the mesh and materials, then applied at every iteration of a solve;
a preconditioner factorized once, then applied every step; a time-step rule with
the step size $\Delta t$ baked in, then marched thousands of times.

Currying gives us a clean way to say exactly this. A constructor function takes
the configuration as its first arguments, and \emph{returns a function} that takes
the run-time data:

\begin{lstlisting}[style=l4style]
-- schematically: configuration first, then the run-time argument
make_operator :: Config -> (Tensor[N] -> Tensor[N])
let apply_A = make_operator cfg in   -- build ONCE, configuration frozen in
let y1 = apply_A x1 in               -- apply MANY times, cheaply
let y2 = apply_A x2 in
...
\end{lstlisting}

\noindent We build \code{apply\_A} once---freezing the configuration \code{cfg}
in---and then apply it again and again to different vectors. This ``build once,
apply many times'' shape is the backbone of \emph{constructed operators}, which
get a full chapter of their own (\cref{ch:operators}). Notice the parentheses in
the signature, \lstinline[style=l4style]|Config -> (Tensor[N] -> Tensor[N])|: we
will explain them carefully in \cref{sec:higherorder}---they are a deliberate
signal that \code{make\_operator} \emph{returns a function}.

\section{Higher-order functions}
\label{sec:higherorder}

We now reach the one genuinely new idea of the chapter---new, that is, to a reader
coming from ordinary imperative programming. It is the idea we have been circling:
a function can \emph{take} another function as an argument, or \emph{return} one as
a result. A function that does either is called a \emph{higher-order function}.

This is the single biggest conceptual lift in the notation, so we build it in two
halves: first a function that \emph{takes} a function, then a function that
\emph{returns} one. Take your time with both; everything in \cref{part:framework}
leans on them.

\subsection{Half one: a function that \emph{takes} a function}

Because a function is a value (\cref{fig:fnvalue}), nothing stops us from writing
a function whose \emph{parameter} is itself a function. Here is the cleanest tiny
example: \code{apply\_twice} takes a function \code{f} and a value \code{x}, and
applies \code{f} to \code{x} twice.

\begin{lstlisting}[style=l4style]
apply_twice :: (Float -> Float) -> Float -> Float
apply_twice f x = f (f x)
\end{lstlisting}

\noindent Read the signature one piece at a time. The first argument has type
\lstinline[style=l4style]|(Float -> Float)|---note the parentheses: this argument
is \emph{a function} from \code{Float} to \code{Float}. The second argument is a
plain \code{Float}. The result is a \code{Float}. The body
\lstinline[style=l4style]|f (f x)| says: apply \code{f} to \code{x}, then apply
\code{f} again to that. We can hand it our \code{square} from before:

\begin{lstlisting}[style=l4style]
apply_twice square 3    -- = square (square 3) = square 9 = 81
add3_twice = apply_twice add3       -- partially applied: freeze f = add3
add3_twice 10                       -- = add3 (add3 10) = add3 13 = 16
\end{lstlisting}

\noindent The parentheses around \lstinline[style=l4style]|(Float -> Float)| in
the signature are essential and worth pausing on. Without them,
\lstinline[style=l4style]|Float -> Float -> Float| would read (by the
right-grouping rule) as a function taking \emph{two} \code{Float}s---an ordinary
two-argument numeric function, like \code{add}. The parentheses say ``no: the
\emph{first} argument is itself a function.'' Grouping is the whole difference
between ``takes a function'' and ``takes two numbers.''

This pattern---pass a function in, let the surrounding function decide how to use
it---is exactly how the combinators of \cref{ch:fold} work. \code{map} takes a
function and applies it to every element of a list. \code{foldl} takes a function
and threads it through a list, accumulating. \code{iterate\_while}, our central
loop combinator, takes a \emph{step} function and applies it over and over. Each
is a function that takes a function; \code{apply\_twice} is the toy version of all
of them.

\begin{figure}[t]
\centering
\begin{tikzpicture}[node distance=10mm]
  % ===== LEFT: a function flowing IN to another function =====
  \begin{scope}
    \node[font=\bfseries\small, simInk] at (0,2.5) {A function flows \emph{in}};
    % the outer machine
    \node[draw=simBlue, fill=simBgBlue, thick, rounded corners=2pt,
          minimum width=30mm, minimum height=20mm] (outer) at (0,0.6) {};
    \node[font=\ttfamily\scriptsize, simInk] at (0,1.35) {apply\_twice};
    % a little function-box being inserted as an argument
    \node[draw=simTeal, fill=simBgGreen, thick, rounded corners=2pt,
          minimum width=13mm, minimum height=7mm] (fin) at (-2.9,0.6)
          {\ttfamily\scriptsize f};
    \draw[flow] (fin) -- (outer.west|-0,0.6);
    \node[font=\scriptsize, simInk] at (0,0.35) {uses \texttt{f} inside};
    \node[draw=simRust,fill=simBgRust,rounded corners=1pt,minimum width=11mm,
          minimum height=7mm] (xin) at (-2.9,-0.6) {\scriptsize $x$};
    \draw[flow] (xin) -- (outer.west|-0,-0.4);
    \draw[flow] (outer.east) -- ++(1.0,0) node[right,font=\scriptsize]{\texttt{f (f x)}};
    \node[align=center,font=\scriptsize\itshape,simGray,text width=42mm]
      at (0,-1.85) {a higher-order function \emph{takes} a function as an argument
      and calls it internally (\code{map}, \code{foldl}, \code{iterate\_while}).};
  \end{scope}
  % ===== RIGHT: a closure flowing OUT, capturing a value =====
  \begin{scope}[shift={(7.4,0)}]
    \node[font=\bfseries\small, simInk] at (0,2.5) {A closure flows \emph{out}};
    % the factory machine
    \node[draw=simBlue, fill=simBgBlue, thick, rounded corners=2pt,
          minimum width=24mm, minimum height=14mm] (fac) at (0,0.9) {};
    \node[font=\ttfamily\scriptsize, simInk] at (0,1.45) {scaled\_by};
    % a captured value going in
    \node[draw=simRust,fill=simBgRust,rounded corners=1pt,minimum width=11mm,
          minimum height=7mm] (cap) at (-2.7,0.9) {\scriptsize $a$};
    \draw[flow] (cap) -- (fac.west);
    \node[font=\scriptsize, simInk] at (0,0.6) {captures $a$};
    % the closure coming out, with the captured a baked in
    \node[draw=simTeal, fill=simBgGreen, thick, rounded corners=2pt,
          minimum width=22mm, minimum height=9mm] (clo) at (0,-1.0) {};
    \node[font=\ttfamily\scriptsize] at (0,-0.75) {\textbackslash v -> a .* v};
    \node[draw=simGold,fill=simBgGold,rounded corners=1pt,font=\scriptsize,
          inner sep=1.3pt] at (0,-1.35) {$a$ baked in};
    \draw[flow] (fac.south) -- (clo.north);
    \node[align=center,font=\scriptsize\itshape,simGray,text width=44mm]
      at (0,-2.25) {a higher-order function \emph{returns} a function that has
      \emph{captured} the value $a$---a \emph{closure}.};
  \end{scope}
\end{tikzpicture}
\caption[The two halves of higher order]{The two ways a function can be
higher-order. \emph{Left:} \code{apply\_twice} \emph{takes} a function \code{f} as
an argument and calls it inside its own body---the pattern of every combinator in
\cref{ch:fold}. \emph{Right:} \code{scaled\_by} (built next) \emph{returns} a
function, and that returned function has \emph{captured} the value $a$ supplied to
the factory---a \emph{closure}, a function with some data baked into it. Functions
flow both in and out.}
\label{fig:higherorder}
\end{figure}

\subsection{Half two: a function that \emph{returns} a function (a closure)}

The other half is a function whose \emph{result} is a function. We have already
half-seen it: \code{add3 = add 3} was a function produced by partially applying
\code{add}. But we can also write a factory \emph{directly}, with the explicit
intention of producing a function. Consider \code{scaled\_by}, which takes a scalar
$a$ and returns the function ``scale a tensor by $a$.''

\begin{lstlisting}[style=l4style]
scaled_by :: Scalar -> (Tensor[N] -> Tensor[N])
scaled_by a = \v -> a .* v
\end{lstlisting}

\noindent Several new pieces here, each small:

\begin{description}[style=nextline,leftmargin=1.5em]
  \item[The codomain \code{(Tensor[N] -> Tensor[N])}.] The
  signature says \code{scaled\_by} takes a \code{Scalar} and returns---the
  parenthesized part---a \emph{function} from \code{Tensor[N]} to \code{Tensor[N]}.
  Those parentheses are the \emph{closure-returning convention} we have been
  flagging: see the box just below.
  \item[The lambda \code{\textbackslash v -> a .* v}.] This is an
  \emph{anonymous function}---a function with no name, written inline. The
  backslash \code{\textbackslash} introduces the parameter (here \code{v}), the
  arrow \code{->} separates the parameter from the body, and
  \lstinline[style=l4style]|a .* v| is the body (the operator \code{.*} is
  scalar-times-tensor; we meet it in \cref{ch:tensors}). Read the whole thing as
  ``the function that, given \code{v}, returns \lstinline[style=l4style]|a .* v|.''
  \item[The capture.] The crucial part: the body
  \lstinline[style=l4style]|a .* v| mentions \code{a}, which is \emph{not} the
  lambda's own parameter---it is the argument that was handed to \code{scaled\_by}.
  The returned function has \emph{captured} \code{a}: it carries that value around
  with it, baked in, ready to use whenever it is finally applied. A function that
  has captured values from where it was created is called a \emph{closure}.
\end{description}

\noindent Now we can use it. Building the closure and applying it are two separate
events, possibly far apart in the program:

\begin{lstlisting}[style=l4style]
let double  = scaled_by 2 in   -- build a closure; 2 captured.  double :: Tensor[N] -> Tensor[N]
let halve   = scaled_by 0.5 in -- a DIFFERENT closure; 0.5 captured
let y = double v in            -- apply later: y = 2 .* v
let z = halve  v in            -- z = 0.5 .* v
(y, z)
\end{lstlisting}

\noindent \code{double} and \code{halve} are two distinct function values, born
from the same factory but carrying different captured scalars. This is the
right-hand cartoon of \cref{fig:higherorder}: a value flows \emph{in} to the
factory, a function flows \emph{out} with that value captured inside it.
\Cref{wex:scaledby} traces this through with concrete numbers.

\begin{notationbox}
\textbf{The closure-returning parenthesis convention.} When a signature's result
(its \emph{codomain}, the type after the last top-level arrow) is itself a
function that you intend to \emph{hold and apply later}, we group it in
parentheses:
\begin{lstlisting}[style=l4style]
scaled_by :: Scalar -> (Tensor[N] -> Tensor[N])   -- "returns a function"
\end{lstlisting}
By the right-grouping rule this is the \emph{same type} as the un-parenthesized
\lstinline[style=l4style]|Scalar -> Tensor[N] -> Tensor[N]|---the parentheses
change no types. They are a \emph{reader-intent marker}: they tell you the
codomain is a closure you will hold and apply, not the last operand of a
multi-argument call. We use them whenever a function's product is itself a
function to be applied later---operator constructors, partial-application
factories. \Cref{ch:operators} introduces a richer spelling,
\lstinline[style=l4style]|Op[in -> out]|, for the special case where the returned
function is an \emph{operator} carrying baked-in parameters; the parenthesis habit
is its plainer sibling.
\end{notationbox}

\subsection{Why ``higher-order'' is worth the effort}

It is fair to ask whether this machinery earns its conceptual cost. It does, and
the reason is the one from \cref{ch:bigidea}: half this book is the discovery
that six different-looking simulators are built from one small set of pieces,
combined in slightly different ways. ``Combined'' is doing real work in that
sentence. The \emph{verbs} of an algorithm (apply this operator, take this inner
product) are ordinary functions; but the \emph{coordination}---do this step, then
this, repeat until converged, map over a family of right-hand sides---is expressed
by functions that \emph{take} the verbs as arguments and orchestrate them. A loop
that works for \emph{any} step function is a higher-order function; a solver that
is configured once and applied many times is a closure. Without functions as
values, we would have to rewrite the coordination from scratch for every verb;
with them, we write the coordination \emph{once} and feed it the verbs. That is
the entire economy of \cref{part:framework}.

\section{Lambda, \code{let}, and the absence of hidden order}

Three loose ends remain, all small, all in service of reading the bodies of
functions fluently.

\subsection{Lambda: anonymous functions}

We met the lambda \lstinline[style=l4style]|\v -> a .* v| above. It is just a
function written without giving it a name. The general form is
\lstinline[style=l4style]|\x -> body|: the backslash, a parameter, an arrow, and a
body expression. We reach for a lambda whenever a function is small and used in
exactly one place---most often as the argument to a higher-order function, where
naming it would be ceremony:

\begin{lstlisting}[style=l4style]
apply_twice (\x -> x + 1) 10   -- = 12; the increment function, written inline
\end{lstlisting}

\noindent Anything written as a lambda could be given a name with \code{let} and a
signature instead; the lambda is a convenience for the throwaway case. The two
forms below are interchangeable:

\begin{lstlisting}[style=l4style]
let inc = \x -> x + 1 in apply_twice inc 10   -- named
apply_twice (\x -> x + 1) 10                   -- inline lambda
\end{lstlisting}

\subsection{\code{let}: naming intermediates}

We have used \code{let} throughout; here is its job stated plainly.
\lstinline[style=l4style]|let name = expr in body| evaluates \code{expr},
binds its value to \code{name}, and then evaluates \code{body} with that name
available. It is how we name an intermediate result so we can refer to it more
than once, or simply to give a complicated computation readable, step-by-step
structure. A chain of \code{let}s reads top to bottom like a recipe:

\begin{lstlisting}[style=l4style]
let r  = b - matvec A x in   -- the residual
let p  = precond r in        -- the preconditioned residual
let a  = dot r p in          -- a scalar
...
\end{lstlisting}

\noindent Each name binds an immutable value; none is ever reassigned. (When a
chain of \code{let}s threads \emph{evolving state}---a counter that advances, a
state vector that steps---we will switch to the \code{do}-block notation of
\cref{ch:monad}, which is the same recipe-like reading with state made explicit.
Plain \code{let} suffices until then.)

\subsection{No hidden order: data dependencies are the only order}

Here is a quietly profound consequence of purity, and the bridge to the rest of
Part~III. Because every function is pure---its output depends only on its
inputs---the \emph{value} of an expression depends on \emph{nothing but the values
of its sub-expressions}. The order in which we happen to write the \code{let}s, or
the order a machine happens to evaluate them, cannot change any answer. The
\emph{only} ordering that matters is forced by the data: if computing \code{p}
needs \code{r}, then \code{r} must be computed first---but two intermediates that
do not depend on each other may be computed in either order, or at the same time.

This means an expression is really a \emph{tree} (more generally, a graph) whose
edges are data dependencies, and that tree carries no order beyond what the edges
demand. \Cref{fig:exprtree} draws one. Read the figure and the lesson of
\cref{ch:bigidea} lands with full force: the \emph{program text is the
dependency graph}, with no hidden arrows from sequencing, because there is no
hidden sequencing---only data flowing along visible edges.

\begin{figure}[t]
\centering
\begin{tikzpicture}[node distance=8mm and 12mm]
  % leaves
  \node[data] (b)  at (-3.2,0)   {\code{b}};
  \node[data] (A)  at (-1.4,-0.9){\code{A}};
  \node[data] (x)  at (0.2,-0.9) {\code{x}};
  % matvec A x
  \node[opbox, minimum width=16mm] (mv) at (-0.6,0.1) {\code{matvec}};
  \draw[flow] (A) -- (mv);
  \draw[flow] (x) -- (mv);
  % r = b - matvec
  \node[opbox, minimum width=12mm] (sub) at (-1.7,1.2) {$-$};
  \draw[flow] (b)  -- (sub);
  \draw[flow] (mv) -- (sub);
  \node[font=\scriptsize\ttfamily, simBlue, left=0.5mm of sub] {r};
  % precond r
  \node[opbox, minimum width=16mm] (pc) at (0.4,1.2) {\code{precond}};
  \draw[flow] (sub) -- (pc);
  \node[font=\scriptsize\ttfamily, simBlue, right=0.5mm of pc] {p};
  % dot r p
  \node[opbox, minimum width=12mm] (dot) at (-0.6,2.5) {\code{dot}};
  \draw[flow] (sub) to[out=70,in=200] (dot);
  \draw[flow] (pc)  -- (dot);
  \node[product, minimum width=10mm, above=7mm of dot] (a) {\code{a}};
  \draw[flow] (dot) -- (a);
  % annotation
  \node[align=left,font=\scriptsize\itshape,simGray,text width=46mm]
    at (4.6,1.2) {The edges are the \emph{only} ordering. \code{precond} must wait
    for \code{r}; but reading \code{b} and forming \code{matvec A x} are unordered
    relative to each other. There is no hidden ``step 1, step 2''---only data
    flowing up the tree.};
\end{tikzpicture}
\caption[An expression is a dependency tree with no hidden order]{The expression
\lstinline[style=l4style]|let r = b - matvec A x in let p = precond r in dot r p|,
drawn as a tree. Leaves are inputs; internal nodes are pure operations; edges are
data dependencies. The only ordering the computation imposes is the one the edges
force (\code{p} after \code{r}, the final \code{dot} after both). Anything the
edges leave unordered may happen in any order, or in parallel---because every node
is a pure function of its inputs and nothing else. The program text \emph{is} this
graph. Reading an algorithm's structure off its written form, exactly here, is the
skill \cref{ch:fold} turns into a method.}
\label{fig:exprtree}
\end{figure}

\section{A first look at operator values}

We close with a forward glance, kept deliberately light. Several times we have met
the pattern ``build a function once with configuration baked in, apply it many
times.'' This pattern is so central---it is how every assembled operator, every
preconditioner, every constructed step works---that the calculus gives it a
dedicated spelling. An \emph{operator value} is written

\begin{lstlisting}[style=l4style]
Op[Tensor[N] -> Tensor[N]]
\end{lstlisting}

\noindent and you should read it as ``a function value from \code{Tensor[N]} to
\code{Tensor[N]} that \emph{also carries baked-in parameters}''---the captured
configuration (matrix entries, a factorization, basis tables) that a plain closure
would hold invisibly, here made a first-class, named thing. An operator
constructor then has a signature like

\begin{lstlisting}[style=l4style]
make_operator :: Config -> Op[Tensor[N] -> Tensor[N]]
\end{lstlisting}

\noindent ``give me a configuration, I return an operator value you apply later.''
The \lstinline[style=l4style]|Op[...]| brackets play exactly the role the
closure-returning parentheses played for \code{scaled\_by}: they announce ``the
result is a thing you hold and apply.'' (Because the brackets already group the
in/out arrow, an \lstinline[style=l4style]|Op[...]| codomain wants no extra outer
parentheses.) Why operators deserve their own spelling rather than just being
plain closures---the answer is about \emph{ownership}, the second of our three
questions, namely \emph{who owns the baked-in parameters} and that they are
read-only for the operator's whole life---is the subject of \cref{ch:operators}.
For now, just recognize the shape: \lstinline[style=l4style]|Op[in -> out]| is a
configured function you apply many times.

\begin{workedexample}{Reading a real signature, token by token}{axpy}
The workhorse update of every linear solver in this book is \code{axpy}---``$a$
times $x$ plus $y$,'' the operation we met in \cref{ch:bigidea}. Here is its
full definition in the notation:
\begin{lstlisting}[style=l4style]
axpy :: Scalar -> Tensor[N] -> Tensor[N] -> Tensor[N]
axpy a x y = a .* x .+ y
\end{lstlisting}
We read \emph{every} token, top line first.
\begin{description}[style=nextline,leftmargin=1.4em,nosep]
  \item[\code{axpy}] the name of the function.
  \item[\code{::}] ``has type.'' Everything after it is the type.
  \item[\code{Scalar}] the type of the \emph{first} argument: a single number $a$.
  \item[the first \code{->}] ``maps to''---and, by right-grouping, ``\dots then
  returns a function that takes\dots''
  \item[\code{Tensor[N]} (second)] the type of the second argument: a length-$N$
  vector $x$. (Shapes are types; \cref{ch:tensors}.)
  \item[\code{Tensor[N]} (third)] the third argument: another length-$N$ vector
  $y$. All three of \code{x}, \code{y}, and the result share the same length $N$.
  \item[\code{Tensor[N]} (last)] the \emph{result} type: a length-$N$ vector. This
  is the codomain---what you get once all three arguments are supplied.
\end{description}
So the signature reads: ``\code{axpy} takes a scalar, then a length-$N$ vector,
then another length-$N$ vector, and returns a fresh length-$N$ vector.'' Now the
body:
\begin{description}[style=nextline,leftmargin=1.4em,nosep]
  \item[\code{axpy a x y}] names the three parameters, in order: \code{a}, \code{x},
  \code{y}.
  \item[\code{a .* x}] the scalar \code{a} times the tensor
  \code{x}, elementwise---a fresh tensor (the \code{.*} dot marks
  scalar-times-tensor; \cref{ch:tensors}).
  \item[\code{.+ y}] add the tensor \code{y} elementwise---a
  fresh tensor, the result.
\end{description}
Three readings worth stating out loud, because each is a habit this chapter is
trying to build. (1) \emph{Nothing is overwritten}: \code{x} and \code{y} name
exactly the same tensors after the call as before; the result is a brand-new
value. (2) \emph{It is curried}: \lstinline[style=l4style]|axpy 2| is already a
legal value---the function ``$2x + y$ awaiting $x$ and $y$''---and
\lstinline[style=l4style]|axpy 2 x| is the function ``$2x + y$ awaiting $y$.''
(3) \emph{It is pure}: hand \code{axpy} the same $a$, $x$, $y$ on any machine, any
day, and you get the identical result; it reads no hidden state and writes none.
\end{workedexample}

\begin{workedexample}{Building and applying a closure, with numbers}{scaledby}
We trace \code{scaled\_by} from \cref{sec:higherorder} through to concrete
numbers, to make the higher-order idea tangible. Recall:
\begin{lstlisting}[style=l4style]
scaled_by :: Scalar -> (Tensor[N] -> Tensor[N])
scaled_by a = \v -> a .* v
\end{lstlisting}
Take $N = 3$ and a concrete input tensor $v = [10, 20, 30]$.

\textbf{Step 1 --- build the closure.} Evaluate \lstinline[style=l4style]|scaled_by 2|.
We supply $a = 2$. By the definition, the result is the lambda
\lstinline[style=l4style]|\v -> a .* v| \emph{with $a$ captured as $2$}:
\begin{lstlisting}[style=l4style]
let double = scaled_by 2 in ...
-- double is the closure (\v -> 2 .* v); the 2 is baked in.
-- double :: Tensor[3] -> Tensor[3]
\end{lstlisting}
No tensor has been touched yet---we have only \emph{produced a function}. This is
the value flowing \emph{out} of the factory in \cref{fig:higherorder} (right).

\textbf{Step 2 --- apply the closure.} Now, possibly much later in the program,
we apply \code{double} to our $v$:
\begin{lstlisting}[style=l4style]
let y = double v in ...   -- substitute v = [10,20,30] into (\v -> 2 .* v)
\end{lstlisting}
Applying the closure substitutes $v = [10,20,30]$ into the captured body
$2 .* v$:
\[
  y \;=\; 2 \cdot [10,\,20,\,30] \;=\; [20,\,40,\,60].
\]

\textbf{Step 3 --- a second, independent closure.} Build another from the same
factory with a different captured scalar:
\begin{lstlisting}[style=l4style]
let halve = scaled_by 0.5 in
let z = halve v in ...     -- z = 0.5 .* [10,20,30] = [5,10,15]
\end{lstlisting}
\code{double} and \code{halve} are two distinct closures, born from one factory,
each carrying its own captured scalar; \code{v} was never disturbed by either
(immutability), and neither closure can interfere with the other (purity). The
two events---\emph{building} the function and \emph{applying} it---are cleanly
separated, which is exactly the ``configure once, apply many times'' shape that
\cref{ch:operators} elevates into constructed operators.
\end{workedexample}

\summarysection
The calculus exists to make three structural questions explicit---what operations
happen, who owns the state, how evolution is coordinated---and this chapter does
the first by teaching the notation for \emph{values and functions}. A
\emph{value} (number, tensor, record, function) is immutable: it never changes
once created; we never overwrite, we compute fresh values and name them with
\code{let}; \code{clone()} is meaningless because nothing can be disturbed. A
\emph{function} is written as a \emph{type signature}---\code{f :: A -> B}, where
\code{::} reads ``has type'' and \code{->} ``maps to''---and a \emph{definition}
applied by juxtaposition; our functions are \emph{pure} (output depends only on
declared inputs, no side effects). A function is itself a \emph{value}, which
gives us two powerful moves. \emph{Currying} reads \code{A -> B -> C} as one
argument at a time, so \emph{partial application} can freeze some arguments now
and supply the rest later (``build once, apply many''). \emph{Higher-order
functions} take a function as an argument (\code{apply\_twice}, the shape of every
combinator) or return one (\code{scaled\_by}, a \emph{closure} that has captured a
value). The closure-returning parenthesis convention
\lstinline[style=l4style]|A -> (B -> C)| marks ``returns a function you hold and
apply later.'' \emph{Lambdas} \lstinline[style=l4style]|\x -> ...| are anonymous
functions; \code{let} names intermediates; and because every function is pure, an
expression is a \emph{dependency tree with no hidden order}---the program text
\emph{is} the dependency graph. The operator-value spelling
\lstinline[style=l4style]|Op[in -> out]|, a configured function carrying baked-in
parameters, previews \cref{ch:operators}.

\furtherreading
The shift from mutable variables to immutable values and first-class functions is
the foundation of typed functional programming; the most readable entry points are
the standard reference for the language whose notation we borrow~\citep{peytonjones2003}
and Wadler's exposition of how a pure language nonetheless models stateful
computation~\citep{wadler1995monads}, which we will need in earnest in
\cref{ch:monad}. No prior background is assumed here or anywhere in the book; each
of the four pillars of the notation gets its own careful chapter---records and
immutability (\cref{ch:records}), tensors and shapes (\cref{ch:tensors}), the fold
and the combinators (\cref{ch:fold}), and the state monad (\cref{ch:monad}).
Readers who want the term ``currying'' in its original setting will find it in any
text on the lambda calculus.

\exercisessection
\begin{enumerate}[nosep]
  \item \textbf{Read this signature.} In one English sentence each, say what these
  functions take and return:
  \code{nrm2 :: Tensor[N] -> Scalar};
  \code{dot :: Tensor[N] -> Tensor[N] -> Scalar};
  \code{matvec :: Tensor[m, n] -> Tensor[n] -> Tensor[m]}.
  For \code{matvec}, why is the result length $m$ rather than $n$?
  \item \textbf{What does it curry to?} Recall \code{axpy} has type
  \code{Scalar -> Tensor[N] -> Tensor[N] -> Tensor[N]}.
  Write the type of \lstinline[style=l4style]|axpy 2|, the type of
  \lstinline[style=l4style]|axpy 2 x| (with \code{x : Tensor[N]}), and the type of
  \lstinline[style=l4style]|axpy 2 x y|. Which of the three is \emph{not} a
  function?
  \item \textbf{Spot the side effect.} One of the following describes a
  \emph{pure} function and one does not: (a) ``returns its first argument plus its
  second''; (b) ``adds its argument to a running total kept in a global counter
  and returns the new total.'' Which is impure, and what specifically about it
  violates purity? Why can the impure one give different answers on the same
  input?
  \item \textbf{Higher-order or not.} For each signature, say whether it is a
  higher-order function and \emph{why}: (a)
  \lstinline[style=l4style]|(Float -> Float) -> Float -> Float|; (b)
  \lstinline[style=l4style]|Float -> Float -> Float|; (c)
  \lstinline[style=l4style]|Scalar -> (Tensor[N] -> Tensor[N])|. For the
  higher-order ones, say whether the function-ness is in an \emph{argument} or in
  the \emph{result}.
  \item \textbf{Build the closure.} Using \code{scaled\_by} from
  \cref{wex:scaledby}, write down (i) the value of
  \lstinline[style=l4style]|scaled_by 3| in words, (ii) its type, and (iii) the
  result of applying it to $[1, 0, 2]$. Then explain why \code{scaled\_by 3} and
  \code{scaled\_by 5} cannot interfere with each other even though they came from
  the same factory.
  \item \textbf{Parentheses matter.} Explain, in terms of the right-grouping rule,
  why \lstinline[style=l4style]|(Float -> Float) -> Float| and
  \lstinline[style=l4style]|Float -> Float -> Float| are \emph{different} types.
  Describe a value of each. (Hint: count how many \emph{numbers} each ultimately
  needs.)
  \item \textbf{No hidden order.} In the dependency tree of \cref{fig:exprtree},
  list every pair of operations whose relative order is \emph{not} forced by the
  data, and one pair whose order \emph{is}. What property of the functions
  guarantees that the unforced pairs may be evaluated in either order without
  changing the answer?
  \item \textbf{The mutable-box trap.} A reader writes: ``after
  \lstinline[style=l4style]|let r = axpy a x y|, the tensor \code{y} now holds
  $a x + y$.'' Identify the exact mental-model error in this sentence and rewrite
  it correctly. What \emph{does} \code{y} refer to after that line?
\end{enumerate}
